% Part: incompleteness
% Chapter: representability-in-q
% Section: prim-rec

\documentclass[../../../include/open-logic-section]{subfiles}

\begin{document}

% SEGMENT OLP-0293-S01

\olfileid{inc}{req}{pri}
\olsection{முதன்மை மீள்வரையறையை உருவகப்படுத்தல்}

இப்போது பீட்டா சார்பைப் பயன்படுத்தி, முதன்மை மீள்வரையறை மூலமான வரையறையை
ஒழுங்கான சிறுமமாக்கலால் ``உருவகப்படுத்த'' முடியும் என்று காட்டலாம்.
$f(\vec x)$, $g(\vec x, y, z)$ ஆகியவை நம்மிடம் உள்ளன என்று கொள்க.
அப்போது $f$,~$g$ ஆகியவற்றிலிருந்து முதன்மை மீள்வரையறையால் வரையறுக்கப்படும்
% SOURCE CORRECTION TA-INC-017: the arguments agree with both defining equations.
$h(\vec x,y)$ என்ற சார்பு பின்வருமாறு அமையும்:
\begin{align*}
h(\vec x, 0) & =  f(\vec x) \\
h(\vec x, y+1) & =  g(\vec x, y, h(\vec x, y)).
\end{align*}
அடிப்படைச் சார்புகளையும், அவற்றிலிருந்து சேர்ப்பு மற்றும் ஒழுங்கான
சிறுமமாக்கலால் வரையறுக்கப்படும் ($\beta$ போன்ற) சார்புகளையும் பயன்படுத்தி,
$f$,~$g$ ஆகியவற்றிலிருந்து சேர்ப்பு மற்றும் ஒழுங்கான சிறுமமாக்கலை மட்டும்
கொண்டு $h$ ஐ வரையறுக்க முடியும் என்று காட்ட வேண்டும்.

\begin{lem}
\ollabel{lem:prim-rec}
$f$,~$g$ ஆகியவற்றிலிருந்து முதன்மை மீள்வரையறையைப் பயன்படுத்தி $h$ ஐ
வரையறுக்க முடியும் எனில், $f$, $g$, சார்புகள் $\Zero$, $\Succ$,
$\Proj{n}{i}$, $\Add$, $\Mult$, $\Char{=}$ ஆகியவற்றிலிருந்து சேர்ப்பு
மற்றும் ஒழுங்கான சிறுமமாக்கலைப் பயன்படுத்தி அதை வரையறுக்க முடியும்.
\end{lem}

\begin{proof}
முதலில், பின்வரும் நிபந்தனைகளை நிறைவுசெய்யும் ஒரு தொடரை~$d$
குறியீடாக்குமாறு அமையும் மீச்சிறு எண்~$d$ ஐத் திருப்பித்தரும் துணைச் சார்பு
$\hat h(\vec x, y)$ ஐ வரையறுக்க:
\begin{enumerate}
\item $(d)_0 = f(\vec x)$, மேலும்
\item ஒவ்வொரு $i < y$ க்கும், $(d)_{i+1} = g(\vec x, i, (d)_i)$,
\end{enumerate}
இங்கு இப்போது $(d)_i$ என்பது $\beta(d,i)$ என்பதன் சுருக்கம்.
வேறுவிதமாகக் கூறினால், $\hat h$ ஆனது
$\tuple{h(\vec x, 0), h(\vec x, 1), \dots, h(\vec x, y)}$ என்ற
தொடரைத் திருப்பித்தருகிறது. $\hat h$ ஐப் பின்வருமாறு எழுதலாம்:
\[
\hat h(\vec x, y) = \umin{d}{(\beta(d,0) = f(\vec x) \land \bforall{i <
  y}{\beta(d,i+1) = g(\vec x, i,\beta(d,i)})}.
\]
கவனிக்க: இங்கு முதன்மை மீள்வரையறை தேவையில்லை; சிறுமமாக்கல் மட்டுமே
தேவை. பீட்டா சார்புத் துணைத்தேற்றம் \olref[bet]{lem:beta} இன்படி, நாம்
சிறுமமாக்கும் சார்பு ஒழுங்கானது.

ஆனால் இப்போது பின்வருவது நம்மிடம் உள்ளது:
\[
h(\vec x, y) = \beta(\hat h(\vec x, y), y),
\]
எனவே அடிப்படைச் சார்புகளிலிருந்து சேர்ப்பு மற்றும் ஒழுங்கான சிறுமமாக்கலை
மட்டும் பயன்படுத்தி $h$ ஐ வரையறுக்க முடியும்.
\end{proof}

\end{document}
